\chapter{Conclusion}
\label{ch:conclusion}

\section{Summary}
This thesis tested whether physics-informed neural networks give accurate and useful tail-risk estimates under
jump-diffusion dynamics, against an exact reference at every step. The answer depends on the setting, and the
study can say where the PINN helps and where it does not.

In one dimension a carefully designed PINN learns the density and its tail to within 1\% in VaR and ES at every
level tested. The design matters more than the network: writing the residual for the log density, resolving the
jump integral with enough quadrature nodes and penalising the normalisation each change the tail error by an order
of magnitude, while tail-adaptive sampling does not help. In one dimension, however, finite differences are faster
and far more accurate, so the 1-D study validates the method rather than motivating it.

The method carries over to several assets without a grid. Collocation, normalisation and risk all use one proposal
built from the equation's coefficients, and the Laplacian can be estimated stochastically. With one architecture and
training budget the ES error stays near 1\% up to five assets, with training time almost independent of the
dimension, in a regime where finite differences are already impractical. At ten assets the method breaks.

An amortised network that takes the volatility level, jump intensity and jump mean as inputs covers a whole range of
market states with one training run and provides risk contributions and scenario curves by differentiation. For the
Merton model, though, Monte Carlo is faster at matched error, because exact simulation is cheap; the PINN does not pay
back its training cost in latency.

On real data the PINN reproduces its jump-diffusion model closely, with daily ES within 0.2--2.8\% of the exact model
and identical breaches. Its backtest performance is therefore that of the model: better than Gaussian, historical
simulation and unfiltered jump-diffusion forecasts, conservative at 97.5\%, and behind the best volatility-filtered
models (filtered historical simulation, DCC, GARCH-$t$, EVT).

The physics is worth a great deal of data. Trained on no samples, the PINN matches a network of identical design
trained on $10^5$ to $10^6$ exact samples; with the few thousand observations available in practice, the data-only
network is 5--60\% off.

\section{Answers to the research questions}
\begin{description}
  \item[RQ1.] With the log-form residual the tail is learned as accurately as the bulk for up to five assets. The tail
    is where the method fails first: in ten dimensions, for a very sharp initial density, and outside the region
    covered by training points.
  \item[RQ2.] Against exact values, ES errors are about 1\% up to five assets. On real data the PINN is a faithful
    surrogate of its model, and calibration is determined by the model, not by the network.
  \item[RQ3.] The equation replaces $10^5$--$10^6$ exact samples.
  \item[RQ4.] The method scales to five assets at roughly constant cost and accuracy and breaks at ten with the
    architecture and budget used here.
  \item[RQ5.] One network answers queries over a parameter box in milliseconds, but for the Merton model Monte Carlo
    is faster at equal error, and there is no break-even.
\end{description}

\section{Practical implications}
\begin{itemize}
  \item A PINN is worth using when the model is expensive to simulate, when a smooth and differentiable density is
    needed (risk contributions, scenario analysis, calibration), or when one model must serve many parameter
    values. For a model as simple as Merton with linear portfolios, Monte Carlo or the closed form is the better tool.
  \item A neural density must not be trusted outside the region where it was trained. Importance-sampling estimates
    should use a trust region, as in the real-data study, or a proposal whose tails are lighter than the network's.
  \item Model risk dominates network error in practice: on real data the choice of model and volatility filter
    mattered more for the backtests than the difference between the PINN and the exact density.
\end{itemize}

\section{Limitations}
\begin{itemize}
  \item All synthetic results use one common jump factor, which keeps the jump integral one-dimensional.
    Idiosyncratic jumps would make it $d$-dimensional and are untested.
  \item Volatility enters only as a daily scale; stochastic volatility, which would double the state dimension, is
    not modelled.
  \item The method was not made to work at ten assets, and the amortised model degrades at long horizons in five.
  \item The initial condition is the exact short-time density; in models without a closed form it would have to be
    approximated, adding an error that this study does not measure.
  \item The real-data evaluation covers two five-asset baskets and one test period (2018--2026).
\end{itemize}

\section{Future work}
\begin{itemize}
  \item \textbf{Higher dimensions:} wider networks and longer training, time-marching (training successive time
    windows), and separable or low-rank architectures, to move the breakdown point beyond ten assets.
  \item \textbf{Richer models:} stochastic volatility (Bates), idiosyncratic jumps and self-exciting (Hawkes) jump
    arrivals. These are also the models where simulation is expensive and the PINN could win on latency.
  \item \textbf{Option-implied calibration:} fitting the jump parameters to option prices, so that the tail reflects
    the market's forward-looking view.
  \item \textbf{Uncertainty quantification:} ensembles or Bayesian PINNs to attach an error bar to each risk estimate,
    which would also detect the far-tail extrapolation failure automatically.
  \item \textbf{Short horizons:} a time rescaling that keeps the initial density resolvable, to remove the
    short-start-time failure.
\end{itemize}
